\section{Security Analysis}
\label{sec:security}

We argue each goal informally, reducing to standard assumptions on the primitives.

\paragraph{G1 against $\mathcal{A}_{\text{dist}}$ and $\mathcal{A}_{\text{prov}}$.}
The adversary sees $C$, the header $h$ and the recipient entries $(enc_i, w_i)$. HPKE in base mode
is IND-CCA2 secure under the Gap-Diffie--Hellman assumption on X25519 and standard assumptions on
HKDF and the AEAD~\citep{rfc9180}, so the $w_i$ reveal nothing about $\cek$ to anyone without some
$sk_{P_i}$. Given that $\cek$ is uniformly random and hidden, confidentiality of $\mathsf{tar}(D)$
follows from the nonce-based AE security of AES-GCM as used by STREAM~\citep{hoang2015online}, since
every chunk nonce is unique under a fresh key. A non-recipient provider gains nothing from
publishing its own keys: it is simply not among the $pk_{P_i}$. It also cannot impersonate a
recipient by relabelling itself, because unwrapping requires the recipient's private key; our test
suite exercises exactly this case.

\paragraph{G2.} Sealing reads only the provider's key document; no step sends a message to the
provider. The owner's confidence that $pk_P$ belongs to $P$ reduces to the WebPKI plus the optional
pinning, multi-vantage and transparency checks. $\mathcal{A}_{\text{dns}}$ can substitute a key
only at the moment of a fetch and only if all these checks fail together; a substituted key is
then recorded in the manifest's \texttt{key\_id} and detectable later. This is the main security
cost of G2 relative to an onboarding handshake, which can pin keys out of band.

\paragraph{G3.} The owner's signature covers the entire manifest, including the payload hash and
length, so under EUF-CMA security of Ed25519 any modification of the manifest or payload is
rejected. STREAM's final-chunk flag rejects truncation even if an adversary could also replace the
hash. The HPKE \emph{info} string contains $\mathsf{SHA256}(h)$, so a recipient entry copied into a
different package (for example, to launder a skill under a new identifier) fails to unwrap: the
recipient's key schedule differs. We verified this transplant case experimentally
(\cref{sec:impl}).

\paragraph{G4 against $\mathcal{A}_{\text{client}}$.} By construction the client holds only a
reference and a licence; neither $\cek$, $sk_P$ nor plaintext ever reaches it. A licence authorises
use through the provider and is useless for decryption. This property depends on the provider's
architecture, not on cryptography: a provider that decrypted in its desktop client would lose it
immediately, for the reasons given in \cref{sec:boundary}.

\paragraph{G5 against $\mathcal{A}_{\text{query}}$.} No cryptographic argument applies: the model
is an oracle that has read the plaintext, and $\mathcal{A}_{\text{query}}$ may query it adaptively.
Output filters are heuristics with a fundamental gap: any filter that blocks all replies carrying
information about the skill also blocks the skill's legitimate use, because applying a rule
reveals something about the rule. Semantic reconstruction over many sessions (paraphrase,
translation into a language the filter does not model, or learning thresholds by probing) cannot
be ruled out. The evaluation in \cref{sec:eval} therefore measures how much each defence reduces
extraction and at what cost to utility.

\paragraph{Revocation.}
Without a handshake there is no revocation list the provider has agreed to honour. The owner's
levers are (i) not issuing new licence tokens, which works with providers that check licences, and
(ii) publishing a new version under a fresh $\cek$ that omits a recipient. Neither retracts what a
former recipient already decrypted. Owners who need enforceable terms, audit feeds or instant
revocation can layer an onboarding agreement or an online key-release service on top of the same
package format.
